% !TeX root = main.tex
\section{工业机器人II}

\subsection{人机协作}

\ref{subsec:intelligent} 小节中，我们通过飞机铆接案例提到，工业机器人和人各有各的长处和短板。传统自动化生产线的做法是把人和机器人彻底隔开。工业4.0描述的人机协作则完全不同，它让人和机器人共用同一个工作空间协同作业，这就要求系统必须能可靠地识别和追踪人的位置，对整个工作空间进行实时监控，一旦有风险就能及时反应。这种方式能支持更灵活的生产流程，不用像传统产线那样改流程就要改设备。

人和机器人在同一空间里干活，紧密程度可以从完全分离到完全合体分成五个等级，如图 \ref{fig:zusammenarbeitsgrade} 所示。

\begin{figure}[htbp]
  \centering
  \includegraphics[width=1\textwidth]{zusammenarbeitsgrade.png}
  \caption{人机协作的五个等级}
  \label{fig:zusammenarbeitsgrade}
\end{figure}

\begin{enumerate}
    \item \textbf{单元（Zelle）}：机器人被围栏围住，和人的工作区完全分离。
    \item \textbf{共存（Koexistenz）}：人和机器人之间没有围栏，但各自的工作范围依然不重叠。
    \item \textbf{同步（Synchronisiert）}：两者的工作范围开始出现重叠，但系统会保证他们不会同时出现在这个重叠区域里。
    \item \textbf{合作（Kooperation）}：人和机器人可以同时出现在共同的工作空间里，但不会共同操作同一个工件。
    \item \textbf{协作（Kollaboration）}：人和机器人共同对同一个工件进行操作。
\end{enumerate}

要实现比较高等级的人机协作，最大的挑战是系统要能把传感、感知和反应整合到一起，实时应对生产过程中各种不确定的变化。要做到这一点需要执行机构、传感器和智能软件三者高度配合协调，还需要引入性能足够强的传感器以及配套的数据处理算法，才能让机器人看得懂周围发生了什么、反应得过来。


\subsection{安全规范与概念}

既然人和机器人要共用一个空间，就必须保证人的安全。分成本质安全设计，从机械结构本身入手，比如让机器人边缘更圆滑、材质更软；另一条是功能安全，靠传感器、控制器这些主动防护功能来实时监测和响应危险。

机械安全标准分三个层级，像金字塔一样，如图 \ref{fig:sicherheitsnormen} 所示。
\begin{itemize}
    \item A类标准讲的是最基础的安全设计原则；
    \item B类标准是某一类通用安全技术要求；
    \item C类标准针对具体某一类机械给出非常具体的安全要求
\end{itemize}
层级越往上越具体、越贴近实际应用。

\begin{figure}[htbp]
  \centering
  \includegraphics[width=0.5\textwidth]{sicherheitsnormen.png}
  \caption{专门为人机协作设计的机器人产品}
  \label{fig:sicherheitsnormen}
\end{figure}

\paragraph{DIN EN ISO 10218：机器人的基本安全标准}

这份标准分两部分。第一部分针对机器人本体，规定了工业机器人和机器人系统特有的危险源，以及对机器人本身（作为系统的一个组成部分）的安全要求。第二部分针对机器人系统与集成，规定了机器人被集成到产线或工作单元时会带来的额外风险，以及在集成安装过程中如何保证安全。

\paragraph{ISO/TS 15066：专门为人机协作补充的规范}

ISO/TS 15066是2016年发布的技术规范，专门针对协作机器人系统的设计。它一方面讲怎么设计协作应用、怎么做危险识别和风险分析，另一方面给出了协作应用需要满足的具体要求，包括安全功能设计、协作工位怎么布置、协作机器人怎么使用和运行。

ISO/TS 15066规定了四种能保证人机安全共处的方法，实际应用中常常组合使用。

方法一是\textbf{安全定向停止（Sicherheitsgerichteter Halt）}：人进入协作空间前，机器人必须先进入安全停止状态。本质上人只能在机器人完全静止时才能靠近它，所以不可能发生碰撞。

方法二是\textbf{手动引导（Handführung）}：同样要求人进入协作空间前机器人先停止，然后激活手动引导模式，这时机器人只能被人用安全的、大幅降低的速度手动操纵。也就是说人始终只能接触静止或被自己缓慢操控的机器人。

方法三是\textbf{速度和间距监控（Geschwindigkeits- \& Abstandsregulierung）}：这是唯一允许机器人和人同时在协作空间内运动的方式。系统会实时监测人和机器人之间的距离，机器人的速度会随着这个距离动态调整。离得越近走得越慢，距离足够小就会停下。这里安全距离的计算遵循（EN ISO 13855）：
\begin{equation}
    S = (K × T) + C
\end{equation}

其中 $S$ 是最小安全距离（mm），$K$ 是人接近的速度参数（mm/s），$T$ 是整个系统的反应加制动时间（s），$C$ 是额外的侵入余量（mm）。这个公式假设机器人是静止的、只有人在移动。

TS 15066在此基础上给出了更精细的动态计算公式，把机器人自身的运动也考虑了进去：
\begin{equation}
    s_p(t₀) = s_h + s_r + s_s + c + z_d + z_r
\end{equation}

其中 \(s_h\) 是人的位移量（mm）、\(s_r\) 是机器人的位移量（mm）、\(s_s\) 是机器人的制动距离（mm）、c 是侵入余量（mm）、\(z_d\) 和 \(z_r\) 分别是人和机器人位置的测量不确定度（mm）。相比前一个公式，这个版本同时考虑了人和机器人双方都在运动、以及定位误差带来的不确定性，算出来的安全距离更贴近人机协作的真实场景。

方法四是\textbf{功率和力限制（Kraft- \& Leistungsreduzierung）}：这是唯一允许人和机器人发生物理接触的方式，靠的是机器人本身在结构设计上就足够温和，再叠加安全功能来降低风险。